\documentclass[11pt]{article}

%============================
% SST house macros (minimal)
%============================
\usepackage{amsmath,amssymb,bm,physics}
\usepackage{geometry}
\geometry{margin=1in}

% Robust swirl arrows
\newcommand{\swirlarrow}{\mathchoice{\mkern-2mu\scriptstyle\boldsymbol{\circlearrowleft}}{\mkern-2mu\scriptstyle\boldsymbol{\circlearrowleft}}{\mkern-2mu\scriptscriptstyle\boldsymbol{\circlearrowleft}}{\mkern-2mu\scriptscriptstyle\boldsymbol{\circlearrowleft}}}
\newcommand{\swirlarrowcw}{\mathchoice{\mkern-2mu\scriptstyle\boldsymbol{\circlearrowright}}{\mkern-2mu\scriptstyle\boldsymbol{\circlearrowright}}{\mkern-2mu\scriptscriptstyle\boldsymbol{\circlearrowright}}{\mkern-2mu\scriptscriptstyle\boldsymbol{\circlearrowright}}}

% Canonical symbols
\newcommand{\vswirl}{\mathbf{v}_{\swirlarrow}}
\newcommand{\vnorm}{\lVert \vswirl \rVert}
\newcommand{\rhoF}{\rho_{\!f}}
\newcommand{\rhoE}{\rho_{\!E}}
\newcommand{\rhoM}{\rho_{\!m}}
\newcommand{\rc}{r_c}

\title{Coarse-Graining Corrections to the SST Mass Functional:\\
A Gradient/Curvature Expansion Inspired by NGM4 Turbulence Closure}
\author{Omar Iskandarani}
\date{Feb 15, 2026}

\begin{document}
\maketitle

\section{Purpose}
This note records a minimal, orthodox-friendly correction to the SST energy/mass functional under coarse-graining at scale $\ell$ (filter width). The structure is inspired by the key fact that filtering quadratic nonlinearities produces a gradient-expansion hierarchy, and that including the $O(\ell^4)$ term can restore physically correct inter-scale energy transfer and stability in LES closures (NGM4).

\section{Setup: Quadratic Coarse-Graining and Subgrid Stress}
Let $\bar{\mathbf{v}} = G_\ell * \mathbf{v}$ be a filtered (coarse-grained) velocity field at scale $\ell$. Define the coarse-grained quadratic mismatch (``subgrid stress'')
\begin{equation}
\tau_{ij} \;\equiv\; \overline{v_i v_j} - \bar v_i \bar v_j.
\end{equation}
For smooth fields and symmetric filters, $\tau_{ij}$ admits a local series in even powers of $\ell$:
\begin{equation}
\tau_{ij} \;\approx\; \kappa_2\,\ell^2\,\partial_k \bar v_i\,\partial_k \bar v_j
\;+\;
\kappa_4\,\ell^4\,\partial_{kl} \bar v_i\,\partial_{kl} \bar v_j
\;+\; O(\ell^6),
\label{eq:tau_expansion}
\end{equation}
with dimensionless constants $\kappa_2,\kappa_4$ determined by the filter (e.g.\ Gaussian). This is the generic SST-agnostic mathematical consequence of filtering quadratic quantities.

\section{Effective energy density with gradient/curvature corrections}
The naive coarse-grained kinetic energy density is
\begin{equation}
\rhoE^{(0)}(\ell) \;=\; \frac12 \rhoF\,\lVert \bar{\mathbf{v}}\rVert^2.
\end{equation}
However, quadratic coarse-graining implies that the energy budget must include SGS contributions that are local surrogates of the nonlocal operator. A minimal local truncation is
\begin{equation}
\rhoE^{\mathrm{eff}}(\ell)
\;=\;
\frac12 \rhoF\Big(
\lVert \bar{\mathbf{v}} \rVert^2
+\kappa_2\,\ell^2\,\lVert \nabla \bar{\mathbf{v}} \rVert^2
+\kappa_4\,\ell^4\,\lVert \nabla\nabla \bar{\mathbf{v}} \rVert^2
\Big),
\label{eq:rhoE_eff}
\end{equation}
where $\lVert \nabla \bar{\mathbf{v}} \rVert^2 = \partial_i \bar v_j\,\partial_i \bar v_j$ and $\lVert \nabla\nabla \bar{\mathbf{v}} \rVert^2 = \partial_{ik}\bar v_j\,\partial_{ik}\bar v_j$.
Dimensional consistency:
\begin{align}
[\rhoF] &= \mathrm{kg\,m^{-3}}, \quad [\bar v] = \mathrm{m\,s^{-1}},\\
[\ell^2 \lVert \nabla \bar{\mathbf{v}}\rVert^2] &= \mathrm{m^2}\cdot (\mathrm{s^{-1}})^2 = \mathrm{m^2\,s^{-2}},\\
\Rightarrow [\rhoE^{\mathrm{eff}}] &= \mathrm{kg\,m^{-1}\,s^{-2}} = \mathrm{J\,m^{-3}}.
\end{align}

\section{Mass functional and resolution dependence}
Define the coarse-grained inertial/rest mass functional by
\begin{equation}
M(\ell) \;\equiv\; \frac{1}{c^2}\int \rhoE^{\mathrm{eff}}(\ell)\,dV.
\label{eq:M_def}
\end{equation}
Then
\begin{equation}
M(\ell) = M_0(\ell) + \Delta M_2(\ell) + \Delta M_4(\ell),
\end{equation}
with
\begin{align}
M_0(\ell) &= \frac{\rhoF}{2c^2}\int \lVert \bar{\mathbf{v}}\rVert^2\,dV,\\
\Delta M_2(\ell) &= \frac{\rhoF\,\kappa_2\,\ell^2}{2c^2}\int \lVert \nabla \bar{\mathbf{v}} \rVert^2\,dV,\\
\Delta M_4(\ell) &= \frac{\rhoF\,\kappa_4\,\ell^4}{2c^2}\int \lVert \nabla\nabla \bar{\mathbf{v}} \rVert^2\,dV.
\end{align}
Scaling estimate for a characteristic variation length $L$:
\begin{equation}
\frac{\Delta M_2}{M_0}\sim \kappa_2\left(\frac{\ell}{L}\right)^2,
\qquad
\frac{\Delta M_4}{M_0}\sim \kappa_4\left(\frac{\ell}{L}\right)^4.
\end{equation}
Hence for localized structures with $L\sim \ell$, gradient/curvature corrections can be $O(1)$ and must be accounted for to obtain an $\ell$-stable mass estimate.

\section{Concrete SST action item: implementation in the SST engine}
\subsection*{Goal}
Implement the two-term corrected estimator
\begin{equation}
M_{(2)}(\ell) \equiv M_0(\ell)+\Delta M_2(\ell),
\end{equation}
and optionally the four-term estimator $M_{(4)}(\ell)$ including $\Delta M_4$. Test sensitivity of your current mass-fit pipeline to $\ell$ and quantify reduction of scale-dependence.

\subsection*{Field inputs}
Assume you already have (from simulation or reconstruction):
\begin{itemize}
\item $\bar{\mathbf{v}}(\mathbf{x})$ on a 3D grid (or a filament/tube parametrization mapped to a grid),
\item grid spacing $\delta x$ and domain volume element $dV$,
\item a chosen coarse-grain scale $\ell$ (often an integer multiple of $\delta x$).
\end{itemize}

\subsection*{Numerical derivatives}
Use spectral derivatives when periodic, or high-order finite differences otherwise. Compute:
\begin{align}
\lVert \bar{\mathbf{v}}\rVert^2 &= \bar v_x^2+\bar v_y^2+\bar v_z^2,\\
\lVert \nabla \bar{\mathbf{v}}\rVert^2 &= \sum_{i\in\{x,y,z\}}\sum_{j\in\{x,y,z\}} (\partial_i \bar v_j)^2,\\
\lVert \nabla\nabla \bar{\mathbf{v}}\rVert^2 &= \sum_{i,k\in\{x,y,z\}}\sum_{j\in\{x,y,z\}} (\partial_{ik}\bar v_j)^2.
\end{align}

\subsection*{Estimator formulas}
\begin{align}
M_0(\ell) &= \frac{\rhoF}{2c^2}\sum_{\text{grid}} \lVert \bar{\mathbf{v}}\rVert^2\,dV,\\
\Delta M_2(\ell) &= \frac{\rhoF\,\kappa_2\,\ell^2}{2c^2}\sum_{\text{grid}} \lVert \nabla \bar{\mathbf{v}}\rVert^2\,dV,\\
\Delta M_4(\ell) &= \frac{\rhoF\,\kappa_4\,\ell^4}{2c^2}\sum_{\text{grid}} \lVert \nabla\nabla \bar{\mathbf{v}}\rVert^2\,dV.
\end{align}
\textbf{Practical default:} start with $\kappa_2=1$ (calibration-free diagnostic), then fit $\kappa_2$ by demanding minimal $\ell$-dependence across a range of $\ell$ values for a fixed configuration.

\section{Appendix: Implementation-ready Python skeleton}
\begin{verbatim}
def mass_functional_with_gradients(vx, vy, vz, dx, rho_f, c,
                                  ell, kappa2=1.0, kappa4=0.0,
                                  derivative_backend="spectral"):
    # vx,vy,vz: 3D arrays of filtered velocity components (bar{v})
    # dx: grid spacing (assume isotropic); dV = dx**3
    # ell: coarse-grain/filter scale in meters
    # kappa2,kappa4: dimensionless coefficients (filter-dependent)

    import numpy as np

    dV = dx**3
    v2 = vx*vx + vy*vy + vz*vz

    # --- compute first derivatives: dvj/dxi ---
    # Replace grad3/hess_norm2 with your engine’s derivative routines
    dvx_dx, dvx_dy, dvx_dz = grad3(vx, dx, backend=derivative_backend)
    dvy_dx, dvy_dy, dvy_dz = grad3(vy, dx, backend=derivative_backend)
    dvz_dx, dvz_dy, dvz_dz = grad3(vz, dx, backend=derivative_backend)

    gradv2 = (dvx_dx**2 + dvx_dy**2 + dvx_dz**2 +
              dvy_dx**2 + dvy_dy**2 + dvy_dz**2 +
              dvz_dx**2 + dvz_dy**2 + dvz_dz**2)

    # --- compute second derivatives if needed ---
    if kappa4 != 0.0:
        d2vx = hess_norm2(vx, dx, backend=derivative_backend)
        d2vy = hess_norm2(vy, dx, backend=derivative_backend)
        d2vz = hess_norm2(vz, dx, backend=derivative_backend)
        hessv2 = d2vx + d2vy + d2vz
    else:
        hessv2 = 0.0

    M0  = (rho_f/(2*c**2)) * np.sum(v2)      * dV
    dM2 = (rho_f*kappa2*ell**2/(2*c**2)) * np.sum(gradv2) * dV
    dM4 = (rho_f*kappa4*ell**4/(2*c**2)) * np.sum(hessv2) * dV

    return {"M0": M0, "dM2": dM2, "dM4": dM4, "M": M0 + dM2 + dM4}
\end{verbatim}

\section{Falsifiable diagnostic}
For a fixed configuration (same physical state), compute $M(\ell)$ over a sweep of $\ell$ (e.g.\ $\ell\in\{1,2,4,8\}\delta x$). The corrected estimator should reduce variation:
\begin{equation}
\mathrm{Var}_\ell\big(M_{(2)}(\ell)\big)\;<\;\mathrm{Var}_\ell\big(M_0(\ell)\big).
\end{equation}
If it does not, the chosen derivative backend, boundary treatment, or assumed locality is insufficient, and a nonlocal (kernel) correction should be considered.

\section{References (minimal)}
\begin{thebibliography}{99}

\bibitem{Jakhar2026NGM4}
K.~Jakhar, Y.~Guan, and P.~Hassanzadeh (2026),
\emph{An Analytical and AI-discovered Stable, Accurate, and Generalizable Subgrid-scale Closure for Geophysical Turbulence},
arXiv:2509.20365v3 (Jan 5, 2026).

\bibitem{Eyink2006}
G.~L. Eyink (2006),
\emph{Multi-scale gradient expansion of the turbulent stress tensor},
Journal of Fluid Mechanics \textbf{549}, 159--190.
DOI: 10.1017/S002211200500759X.

\bibitem{Sagaut2006}
P.~Sagaut (2006),
\emph{Large Eddy Simulation for Incompressible Flows: An Introduction},
Springer.
DOI: 10.1007/978-3-540-26347-2.

\end{thebibliography}

\end{document}
